\begin{titlepage}
	\definecolor{purplemain}{RGB}{66,44,118}     % 深紫（文字与线条用色）
	\definecolor{purplemid}{RGB}{166,144,210}    % 中紫
	\definecolor{purplelight}{RGB}{228,221,247}  % 浅紫
	\begin{tikzpicture}[remember picture, overlay]
		% ---- 1. 浅紫渐变铺满整个封面 ----
		\shade[top color=purplelight, middle color=purplemid!72, bottom color=purplemid]
			(current page.south west) rectangle (current page.north east);

		% ---- 2. 细点阵（模拟网格，微弱纹理）----
		\foreach \x in {1,...,20} {
			\foreach \y in {1,...,29} {
				\fill[purplemain, opacity=0.05] (\x,\y) circle[radius=0.018cm];
			}
		}

		% ---- 3. 主标题后方的柔光晕 ----
		\fill[white, opacity=0.28] ([yshift=-8.2cm]current page.north) circle[radius=7.0cm];
		\fill[white, opacity=0.20] ([yshift=-8.2cm]current page.north) circle[radius=4.6cm];

		% ---- 4. 右上：闭环反馈控制方块图（控制器 + 被控对象 + 反馈）----
		\begin{scope}[shift={([xshift=-6.2cm,yshift=-4.7cm]current page.north east)}]
			% 输入 r(s)
			\node[anchor=east, text=purplemain, font=\footnotesize\itshape] at (0.02,1.7) {$r(s)$};
			% 求和点 ⊕
			\draw[purplemain, line width=0.9pt] (0.85,1.7) circle[radius=0.12cm];
			\draw[purplemain, line width=0.7pt] (0.78,1.7) -- (0.92,1.7);
			\draw[purplemain, line width=0.7pt] (0.85,1.63) -- (0.85,1.77);
			% r -> ⊕
			\draw[purplemain, line width=0.8pt, -{Stealth}] (0.30,1.7) -- (0.73,1.7);
			% ⊕ -> Gc
			\draw[purplemain, line width=0.8pt, -{Stealth}] (0.97,1.7) -- (1.50,1.7);
			% 控制器方块 Gc
			\draw[purplemain, line width=1.0pt, fill=white, fill opacity=0.85] (1.50,1.25) rectangle (2.60,2.15);
			\node[anchor=center, text=purplemain, font=\footnotesize] at (2.05,1.7) {$G_c(s)$};
			% Gc -> Gp
			\draw[purplemain, line width=0.8pt, -{Stealth}] (2.60,1.7) -- (3.00,1.7);
			% 被控对象方块 Gp
			\draw[purplemain, line width=1.0pt, fill=white, fill opacity=0.85] (3.00,1.25) rectangle (4.10,2.15);
			\node[anchor=center, text=purplemain, font=\footnotesize] at (3.55,1.7) {$G_p(s)$};
			% Gp -> y
			\draw[purplemain, line width=0.8pt, -{Stealth}] (4.10,1.7) -- (4.48,1.7);
			\node[anchor=west, text=purplemain, font=\footnotesize\itshape] at (4.50,1.7) {$y(s)$};
			% 引出反馈支路
			\draw[purplemain, line width=0.8pt] (4.25,1.7) -- (4.25,0.6);
			\draw[purplemain, line width=0.8pt, -{Stealth}] (4.25,0.6) -- (4.10,0.6);
			% 反馈方块 H
			\draw[purplemain, line width=1.0pt, fill=white, fill opacity=0.85] (3.00,0.25) rectangle (4.10,0.95);
			\node[anchor=center, text=purplemain, font=\footnotesize] at (3.55,0.6) {$H(s)$};
			% H -> ⊕（负反馈回到求和点）
			\draw[purplemain, line width=0.8pt, -{Stealth}] (3.00,0.6) -- (0.85,0.6) -- (0.85,1.58);
		\end{scope}

		% ---- 5. 左下：复平面极点-零点图（s 平面）----
		\begin{scope}[shift={([xshift=0.9cm,yshift=0.9cm]current page.south west)}]
			% 坐标轴：实轴（σ）与虚轴（jω）
			\draw[purplemain, opacity=0.85, line width=0.9pt, -{Stealth}] (0.2,1.15) -- (2.9,1.15);
			\draw[purplemain, opacity=0.85, line width=0.9pt, -{Stealth}] (1.4,0.10) -- (1.4,2.30);
			\node[anchor=north, text=purplemain, opacity=0.85, font=\footnotesize\itshape] at (2.9,1.15) {$\sigma$};
			\node[anchor=east, text=purplemain, opacity=0.85, font=\footnotesize\itshape] at (1.4,2.36) {$j\omega$};
			% 极点 ×（左半平面，共轭对）
			\draw[purplemain, line width=1.1pt] (0.75,1.65) ++(-0.09,-0.09) -- ++(0.18,0.18);
			\draw[purplemain, line width=1.1pt] (0.75,1.65) ++(-0.09,0.09) -- ++(0.18,-0.18);
			\draw[purplemain, line width=1.1pt] (0.75,0.65) ++(-0.09,-0.09) -- ++(0.18,0.18);
			\draw[purplemain, line width=1.1pt] (0.75,0.65) ++(-0.09,0.09) -- ++(0.18,-0.18);
			% 零点 ○（实轴上）
			\draw[purplemain, line width=1.1pt] (2.05,1.15) circle[radius=0.09cm];
			% 图注
			\node[anchor=west, text=purplemain, opacity=0.85, font=\footnotesize\itshape] at (1.55,1.62) {极点};
			\node[anchor=west, text=purplemain, opacity=0.85, font=\footnotesize\itshape] at (2.24,1.00) {零点};
		\end{scope}

		% ---- 6. 金色细边框 ----
		\draw[gold, line width=0.7pt, opacity=0.6]
			([shift={(0.85,0.85)}]current page.south west) rectangle
			([shift={(-0.85,-0.85)}]current page.north east);

		% ---- 7. 顶部眉题（英文小字 + 金色短线）----
		\coordinate (T) at ([yshift=-3.5cm]current page.north);
		\draw[gold, line width=1.0pt] (T) ++(-3.4cm,0) -- ++(1.0cm,0);
		\draw[gold, line width=1.0pt] (T) ++(2.4cm,0) -- ++(1.0cm,0);
		\node[anchor=north, text=purplemain, opacity=0.85, align=center] at (T) {
			{\sffamily\fontsize{13}{16}\selectfont
				\uppercase{Control \ Theory}}
		};

		% ---- 8. 主标题（居中，使用 \notename）----
		\node[anchor=center, align=center, text=purplemain] at ([yshift=-8.0cm]current page.north) {%
			{\fontsize{40}{50}\sffamily\bfseries \notename}
		};

		% ---- 9. 金色双细线 + 极点纹饰（× 极点符号）----
		\coordinate (C) at ([yshift=-11.4cm]current page.north);
		\draw[gold, line width=1.1pt] (C) ++(-2.2cm,0) -- ++(1.55cm,0);
		\draw[gold, line width=1.1pt] (C) ++(0.65cm,0) -- ++(1.55cm,0);
		\draw[gold, line width=0.9pt] (C) circle[radius=0.18cm];
		\draw[gold, line width=0.8pt] (C) ++(-0.09cm,-0.09cm) -- ++(0.18cm,0.18cm);
		\draw[gold, line width=0.8pt] (C) ++(-0.09cm,0.09cm) -- ++(0.18cm,-0.18cm);

		% ---- 10. 副标题 ----
		\node[anchor=north, text=purplemain, opacity=0.9, align=center]
			at ([yshift=-12.2cm]current page.north) {
			{\fontsize{16}{20}\sffamily 学\ \ 习\ \ 笔\ \ 记}
		};

		% ---- 11. 底部作者、日期信息 ----
		\coordinate (B) at ([yshift=3.4cm]current page.south);
		\draw[gold, line width=0.8pt, opacity=0.7] (B) ++(-1.1cm,0.85cm) -- ++(2.2cm,0);
		\node[anchor=south, text=purplemain, align=center] at (B) {
			\begin{tabular}{c}
				{\fontsize{20}{24}\sffamily\bfseries \noteauthor} \\[0.4cm]
				{\large \sffamily \today}
			\end{tabular}
		};
	\end{tikzpicture}
\end{titlepage}
